%%%%%%%%%%%%%%%%%%%%%%%%%%
%%% 08-samuel/2-13.tex %%%
%%%%%%%%%%%%%%%%%%%%%%%%%%

\Chapter{13}

\Verse{1}
{
	\hJ{וַיְהִי אַחֲרֵי־כֵן וּלְאַבְשָׁלוֹם בֶּן־דָּוִד אָחוֹת יָפָה וּשְׁמָהּ תָּמָר וַיֶּאֱהָבֶהָ אַמְנוֹן בֶּן־דָּוִד׃}
}
{
	\eJ{
		This happened sometime afterward: Absalom son of David had a
		beautiful sister named Tamar, and Amnon son of David became
		infatuated with her.
	}
}
{
}

\Verse{2}
{
	\hJ{וַיֵּצֶר לְאַמְנוֹן לְהִתְחַלּוֹת בַּעֲבוּר תָּמָר אֲחֹתוֹ כִּי בְתוּלָה הִיא וַיִּפָּלֵא בְּעֵינֵי אַמְנוֹן לַעֲשׂוֹת לָהּ מְאוּמָה׃}
}
{
	\eJ{
		Amnon was so distraught because of his [half-]sister Tamar that he
		became sick; for she was a virgin, and it seemed impossible to
		Amnon to do anything to her.
	}
}
{
}

\Verse{3}
{
	\hJ{וּלְאַמְנוֹן רֵעַ וּשְׁמוֹ יוֹנָדָב בֶּן־שִׁמְעָה אֲחִי דָוִד וְיוֹנָדָב אִישׁ חָכָם מְאֹד׃}
}
{
	\eJ{
		Amnon had a friend named Jonadab, the son of David’s brother
		Shimah; Jonadab was a very clever man.
	}
}
{
}

\Verse{4}
{
	\hJ{וַיֹּאמֶר לוֹ מַדּוּעַ אַתָּה כָּכָה דַּל בֶּן־הַמֶּלֶךְ בַּבֹּקֶר בַּבֹּקֶר הֲלוֹא תַּגִּיד לִי וַיֹּאמֶר לוֹ אַמְנוֹן אֶת־תָּמָר אֲחוֹת אַבְשָׁלֹם אָחִי אֲנִי אֹהֵב׃}
}
{
	\eJ{
		He asked him, “Why are you so dejected, O prince, morning after
		morning? Tell me!” Amnon replied, “I am in love with Tamar, the
		sister of my brother Absalom!”
	}
}
{
}

\Verse{5}
{
	\hJ{
		וַיֹּאמֶר לוֹ יְהוֹנָדָב שְׁכַב עַל־מִשְׁכָּבְךָ וְהִתְחָל וּבָא אָבִיךָ לִרְאוֹתֶךָ וְאָמַרְתָּ אֵלָיו תָּבֹא נָא תָמָר אֲחוֹתִי וְתַבְרֵנִי לֶחֶם וְעָשְׂתָה לְעֵינַי
		אֶת־הַבִּרְיָה לְמַעַן אֲשֶׁר אֶרְאֶה וְאָכַלְתִּי מִיָּדָהּ׃
	}
}
{
	\eJ{
		Jonadab said to him, “Lie down in your bed and pretend you are
		sick. When your father comes to see you, say to him, ‘Let my
		sister Tamar come and give me something to eat. Let her prepare
		the food in front of me, so that I may look on, and let her serve
		it to me.’”
	}
}
{
}

\Verse{6}
{
	\hJ{וַיִּשְׁכַּב אַמְנוֹן וַיִּתְחָל וַיָּבֹא הַמֶּלֶךְ לִרְאוֹתוֹ וַיֹּאמֶר אַמְנוֹן אֶל־הַמֶּלֶךְ תָּבוֹא־נָא תָּמָר אֲחֹתִי וּתְלַבֵּב לְעֵינַי שְׁתֵּי לְבִבוֹת וְאֶבְרֶה מִיָּדָהּ׃}
}
{
	\eJ{
		Amnon lay down and pretended to be sick. The king came to see him,
		and Amnon said to the king, “Let my sister Tamar come and prepare
		a couple of cakes in front of me, and let her bring them to me.”
	}
}
{
}

\Verse{7}
{
	\hJ{וַיִּשְׁלַח דָּוִד אֶל־תָּמָר הַבַּיְתָה לֵאמֹר לְכִי נָא בֵּית אַמְנוֹן אָחִיךְ וַעֲשִׂי־לוֹ הַבִּרְיָה׃}
}
{
	\eJ{
		David sent a message to Tamar in the palace, “Please go to the
		house of your brother Amnon and prepare some food for him.”
	}
}
{
}

\Verse{8}
{
	\hJ{וַתֵּלֶךְ תָּמָר בֵּית אַמְנוֹן אָחִיהָ וְהוּא שֹׁכֵב וַתִּקַּח אֶת־הַבָּצֵק (ותלוש) [וַתָּלשׁ] וַתְּלַבֵּב לְעֵינָיו וַתְּבַשֵּׁל אֶת־הַלְּבִבוֹת׃}
}
{
	\eJ{
		Tamar went to the house of her brother Amnon, who was in bed. She
		took dough and kneaded it into cakes in front of him, and cooked
		the cakes.
	}
}
{
	Qere Ketiv.
}

\Verse{9}
{
	\hJ{וַתִּקַּח אֶת־הַמַּשְׂרֵת וַתִּצֹק לְפָנָיו וַיְמָאֵן לֶאֱכוֹל וַיֹּאמֶר אַמְנוֹן הוֹצִיאוּ כל־אִישׁ מֵעָלַי וַיֵּצְאוּ כל־אִישׁ מֵעָלָיו׃}
}
{
	\eJ{
		She took the pan and set out [the cakes], but Amnon refused to eat
		and ordered everyone to withdraw. After everyone had withdrawn,
	}
}
{
}

\Verse{10}
{
	\hJ{וַיֹּאמֶר אַמְנוֹן אֶל־תָּמָר הָבִיאִי הַבִּרְיָה הַחֶדֶר וְאֶבְרֶה מִיָּדֵךְ וַתִּקַּח תָּמָר אֶת־הַלְּבִבוֹת אֲשֶׁר עָשָׂתָה וַתָּבֵא לְאַמְנוֹן אָחִיהָ הֶחָדְרָה׃}
}
{
	\eJ{
		Amnon said to Tamar, “Bring the food inside and feed me.” Tamar
		took the cakes she had made and brought them to her brother
		inside.
	}
}
{
}

\Verse{11}
{
	\hJ{וַתַּגֵּשׁ אֵלָיו לֶאֱכֹל וַיַּחֲזֶק־בָּהּ וַיֹּאמֶר לָהּ בּוֹאִי שִׁכְבִי עִמִּי אֲחוֹתִי׃}
}
{
	\eJ{
		But when she served them to him, he caught hold of her and said to
		her, “Come lie with me, sister.”
	}
}
{
}

\Verse{12}
{
	\hJ{וַתֹּאמֶר לוֹ אַל־אָחִי אַל־תְּעַנֵּנִי כִּי לֹא־יֵעָשֶׂה כֵן בְּיִשְׂרָאֵל אַל־תַּעֲשֵׂה אֶת־הַנְּבָלָה הַזֹּאת׃}
}
{
	\eJ{
		But she said to him, “Don’t, brother. Don’t force me. Such things
		are not done in Israel! Don’t do such a vile thing!
	}
}
{
}

\Verse{13}
{
	\hJ{וַאֲנִי אָנָה אוֹלִיךְ אֶת־חֶרְפָּתִי וְאַתָּה תִּהְיֶה כְּאַחַד הַנְּבָלִים בְּיִשְׂרָאֵל וְעַתָּה דַּבֶּר־נָא אֶל־הַמֶּלֶךְ כִּי לֹא יִמְנָעֵנִי מִמֶּךָּ׃}
}
{
	\eJ{
		Where will I carry my shame? And you, you will be like any of the
		scoundrels in Israel! Please, speak to the king; he will not
		refuse me to you.”
	}
}
{
}

\Verse{14}
{
	\hJ{וְלֹא אָבָה לִשְׁמֹעַ בְּקוֹלָהּ וַיֶּחֱזַק מִמֶּנָּה וַיְעַנֶּהָ וַיִּשְׁכַּב אֹתָהּ׃}
}
{
	\eJ{
		But he would not listen to her; he overpowered her and lay with
		her by force.
	}
}
{
}

\Verse{15}
{
	\hJ{וַיִּשְׂנָאֶהָ אַמְנוֹן שִׂנְאָה גְּדוֹלָה מְאֹד כִּי גְדוֹלָה הַשִּׂנְאָה אֲשֶׁר שְׂנֵאָהּ מֵאַהֲבָה אֲשֶׁר אֲהֵבָהּ וַיֹּאמֶר־לָהּ אַמְנוֹן קוּמִי לֵכִי׃}
}
{
	\eJ{
		Then Amnon felt a very great loathing for her; indeed, his
		loathing for her was greater than the passion he had felt for her.
		And Amnon said to her, “Get out!”
	}
}
{
}

\Verse{16}
{
	\hJ{וַתֹּאמֶר לוֹ אַל־אוֹדֹת הָרָעָה הַגְּדוֹלָה הַזֹּאת מֵאַחֶרֶת אֲשֶׁר־עָשִׂיתָ עִמִּי לְשַׁלְּחֵנִי וְלֹא אָבָה לִשְׁמֹעַ לָהּ׃}
}
{
	\eJ{
		She pleaded with him, “Please don’t commit this wrong; to send me
		away would be even worse than the first wrong you committed
		against me.” But he would not listen to her.
	}
}
{
}

\Verse{17}
{
	\hJ{וַיִּקְרָא אֶת־נַעֲרוֹ מְשָׁרְתוֹ וַיֹּאמֶר שִׁלְחוּ־נָא אֶת־זֹאת מֵעָלַי הַחוּצָה וּנְעֹל הַדֶּלֶת אַחֲרֶיהָ׃}
}
{
	\eJ{
		He summoned his young attendant and said, “Get that woman out of
		my presence, and bar the door behind her.”—
	}
}
{
}

\Verse{18}
{
	\hJ{וְעָלֶיהָ כְּתֹנֶת פַּסִּים כִּי כֵן תִּלְבַּשְׁןָ בְנוֹת־הַמֶּלֶךְ הַבְּתוּלֹת מְעִילִים וַיֹּצֵא אוֹתָהּ מְשָׁרְתוֹ הַחוּץ וְנָעַל הַדֶּלֶת אַחֲרֶיהָ׃}
}
{
	\eJ{
		She was wearing an ornamented tunic, for maiden princesses were
		customarily dressed in such garments. —His attendant took her
		outside and barred the door after her.
	}
}
{
}

\Verse{19}
{
	\hJ{וַתִּקַּח תָּמָר אֵפֶר עַל־רֹאשָׁהּ וּכְתֹנֶת הַפַּסִּים אֲשֶׁר עָלֶיהָ קָרָעָה וַתָּשֶׂם יָדָהּ עַל־רֹאשָׁהּ וַתֵּלֶךְ הָלוֹךְ וְזָעָקָה׃}
}
{
	\eJ{
		Tamar put dust on her head and rent the ornamented tunic she was
		wearing; she put her hands on her head, and walked away, screaming
		loudly as she went.
	}
}
{
}

\Verse{20}
{
	\hJ{
		וַיֹּאמֶר אֵלֶיהָ אַבְשָׁלוֹם אָחִיהָ הַאֲמִינוֹן אָחִיךְ הָיָה עִמָּךְ וְעַתָּה אֲחוֹתִי הַחֲרִישִׁי אָחִיךְ הוּא אַל־תָּשִׁיתִי אֶת־לִבֵּךְ לַדָּבָר הַזֶּה וַתֵּשֶׁב תָּמָר
		וְשֹׁמֵמָה בֵּית אַבְשָׁלוֹם אָחִיהָ׃
	}
}
{
	\eJ{
		Her brother Absalom said to her, “Was it your brother Amnon who
		did this to you? For the present, sister, keep quiet about it; he
		is your brother. Don’t brood over the matter.” And Tamar remained
		in her brother Absalom’s house, forlorn.
	}
}
{
}

\Verse{21}
{
	\hJ{וְהַמֶּלֶךְ דָּוִד שָׁמַע אֵת כּל־הַדְּבָרִים הָאֵלֶּה וַיִּחַר לוֹ מְאֹד׃}
}
{
	\eJ{
		When King David heard about all this, he was greatly upset.
	}
}
{
}

\Verse{22}
{
	\hJ{וְלֹא־דִבֶּר אַבְשָׁלוֹם עִם־אַמְנוֹן לְמֵרָע וְעַד־טוֹב כִּי־שָׂנֵא אַבְשָׁלוֹם אֶת־אַמְנוֹן עַל־דְּבַר אֲשֶׁר עִנָּה אֵת תָּמָר אֲחֹתוֹ׃}
}
{
	\eJ{
		Absalom didn’t utter a word to Amnon, good or bad; but Absalom
		hated Amnon because he had violated his sister Tamar.
	}
}
{
}

\Verse{23}
{
	\hJ{וַיְהִי לִשְׁנָתַיִם יָמִים וַיִּהְיוּ גֹזְזִים לְאַבְשָׁלוֹם בְּבַעַל חָצוֹר אֲשֶׁר עִם־אֶפְרָיִם וַיִּקְרָא אַבְשָׁלוֹם לְכל־בְּנֵי הַמֶּלֶךְ׃}
}
{
	\eJ{
		Two years later, when Absalom was having his flocks sheared at
		Baal-hazor near Ephraim, Absalom invited all the king’s sons.
	}
}
{
}

\Verse{24}
{
	\hJ{וַיָּבֹא אַבְשָׁלוֹם אֶל־הַמֶּלֶךְ וַיֹּאמֶר הִנֵּה־נָא גֹזְזִים לְעַבְדֶּךָ יֵלֶךְ־נָא הַמֶּלֶךְ וַעֲבָדָיו עִם־עַבְדֶּךָ׃}
}
{
	\eJ{
		And Absalom came to the king and said, “Your servant is having his
		flocks sheared. Would Your Majesty and your retinue accompany your
		servant?”
	}
}
{
}

\Verse{25}
{
	\hJ{וַיֹּאמֶר הַמֶּלֶךְ אֶל־אַבְשָׁלוֹם אַל־בְּנִי אַל־נָא נֵלֵךְ כֻּלָּנוּ וְלֹא נִכְבַּד עָלֶיךָ וַיִּפְרץ־בּוֹ וְלֹא־אָבָה לָלֶכֶת וַיְבָרְכֵהוּ׃}
}
{
	\eJ{
		But the king answered Absalom, “No, my son. We must not all come,
		or we’ll be a burden to you.” He urged him, but he would not go,
		and he said good-bye to him.
	}
}
{
}

\Verse{26}
{
	\hJ{וַיֹּאמֶר אַבְשָׁלוֹם וָלֹא יֵלֶךְ־נָא אִתָּנוּ אַמְנוֹן אָחִי וַיֹּאמֶר לוֹ הַמֶּלֶךְ לָמָּה יֵלֵךְ עִמָּךְ׃}
}
{
	\eJ{
		Thereupon Absalom said, “In that case, let my brother Amnon come
		with us,” to which the king replied, “He shall not go with you.”
	}
}
{
}

\Verse{27}
{
	\hJ{וַיִּפְרץ־בּוֹ אַבְשָׁלוֹם וַיִּשְׁלַח אִתּוֹ אֶת־אַמְנוֹן וְאֵת כּל־בְּנֵי הַמֶּלֶךְ׃}
}
{
	\eJ{
		But Absalom urged him, and he sent with him Amnon and all the
		other princes.
	}
}
{
}

\Verse{28}
{
	\hJ{
		וַיְצַו אַבְשָׁלוֹם אֶת־נְעָרָיו לֵאמֹר רְאוּ נָא כְּטוֹב לֵב־אַמְנוֹן בַּיַּיִן וְאָמַרְתִּי אֲלֵיכֶם הַכּוּ אֶת־אַמְנוֹן וַהֲמִתֶּם אֹתוֹ אַל־תִּירָאוּ הֲלוֹא כִּי אָנֹכִי
		צִוִּיתִי אֶתְכֶם חִזְקוּ וִהְיוּ לִבְנֵי־חָיִל׃
	}
}
{
	\eJ{
		Now Absalom gave his attendants these orders: “Watch, and when
		Amnon is merry with wine and I tell you to strike down Amnon, kill
		him! Don’t be afraid, for it is I who give you the order. Act with
		determination, like brave men!”
	}
}
{
}

\Verse{29}
{
	\hJ{וַיַּעֲשׂוּ נַעֲרֵי אַבְשָׁלוֹם לְאַמְנוֹן כַּאֲשֶׁר צִוָּה אַבְשָׁלוֹם וַיָּקֻמוּ כּל־בְּנֵי הַמֶּלֶךְ וַיִּרְכְּבוּ אִישׁ עַל־פִּרְדּוֹ וַיָּנֻסוּ׃}
}
{
	\eJ{
		Absalom’s attendants did to Amnon as Absalom had ordered;
		whereupon all the other princes mounted their mules and fled.
	}
}
{
}

\Verse{30}
{
	\hJ{וַיְהִי הֵמָּה בַדֶּרֶךְ וְהַשְּׁמֻעָה בָאָה אֶל־דָּוִד לֵאמֹר הִכָּה אַבְשָׁלוֹם אֶת־כּל־בְּנֵי הַמֶּלֶךְ וְלֹא־נוֹתַר מֵהֶם אֶחָד׃}
}
{
	\eJ{
		They were still on the road when a rumor reached David that
		Absalom had killed all the princes, and that not one of them had
		survived.
	}
}
{
}

\Verse{31}
{
	\hJ{וַיָּקם הַמֶּלֶךְ וַיִּקְרַע אֶת־בְּגָדָיו וַיִּשְׁכַּב אָרְצָה וְכל־עֲבָדָיו נִצָּבִים קְרֻעֵי בְגָדִים׃}
}
{
	\eJ{
		At this, David rent his garment and lay down on the ground, and
		all his courtiers stood by with their clothes rent.
	}
}
{
}

\Verse{32}
{
	\hJ{
		וַיַּעַן יוֹנָדָב בֶּן־שִׁמְעָה אֲחִי־דָוִד וַיֹּאמֶר אַל־יֹאמַר אֲדֹנִי אֵת כּל־הַנְּעָרִים בְּנֵי־הַמֶּלֶךְ הֵמִיתוּ כִּי־אַמְנוֹן לְבַדּוֹ מֵת כִּי־עַל־פִּי אַבְשָׁלוֹם הָיְתָה
		שׂוּמָה מִיּוֹם עַנֹּתוֹ אֵת תָּמָר אֲחֹתוֹ׃
	}
}
{
	\eJ{
		But Jonadab, the son of David’s brother Shimah, said, “My lord
		must not think that all the young princes have been killed. Only
		Amnon is dead; for this has been decided by Absalom ever since his
		sister Tamar was violated.
	}
}
{
}

\Verse{33}
{
	\hJ{וְעַתָּה אַל־יָשֵׂם אֲדֹנִי הַמֶּלֶךְ אֶל־לִבּוֹ דָּבָר לֵאמֹר כּל־בְּנֵי הַמֶּלֶךְ מֵתוּ כִּי־(אם)־אַמְנוֹן לְבַדּוֹ מֵת׃}
}
{
	\eJ{
		So my lord the king must not think for a moment that all the
		princes are dead; Amnon alone is dead.”
	}
}
{
}

\Verse{34}
{
	\hJ{וַיִּבְרַח אַבְשָׁלוֹם וַיִּשָּׂא הַנַּעַר הַצֹּפֶה אֶת־עֵינָו וַיַּרְא וְהִנֵּה עַם־רַב הֹלְכִים מִדֶּרֶךְ אַחֲרָיו מִצַּד הָהָר׃}
}
{
	\eJ{
		Meanwhile Absalom had fled. The watchman on duty looked up and saw
		a large crowd coming from the road to his rear, from the side of
		the hill.
	}
}
{
}

\Verse{35}
{
	\hJ{וַיֹּאמֶר יוֹנָדָב אֶל־הַמֶּלֶךְ הִנֵּה בְנֵי־הַמֶּלֶךְ בָּאוּ כִּדְבַר עַבְדְּךָ כֵּן הָיָה׃}
}
{
	\eJ{
		Jonadab said to the king, “See, the princes have come! It is just
		as your servant said.”
	}
}
{
}

\Verse{36}
{
	\hJ{וַיְהִי כְּכַלֹּתוֹ לְדַבֵּר וְהִנֵּה בְנֵי־הַמֶּלֶךְ בָּאוּ וַיִּשְׂאוּ קוֹלָם וַיִּבְכּוּ וְגַם־הַמֶּלֶךְ וְכל־עֲבָדָיו בָּכוּ בְּכִי גָּדוֹל מְאֹד׃}
}
{
	\eJ{
		As he finished speaking, the princes came in and broke into
		weeping; and David and all his courtiers wept bitterly, too.
	}
}
{
}

\Verse{37}
{
	\hJ{וְאַבְשָׁלוֹם בָּרַח וַיֵּלֶךְ אֶל־תַּלְמַי בֶּן־[עַמִּיהוּד] (עמיחור) מֶלֶךְ גְּשׁוּר וַיִּתְאַבֵּל עַל־בְּנוֹ כּל־הַיָּמִים׃}
}
{
	\eJ{
		Absalom had fled, and he came to Talmai son of Ammihud, king of
		Geshur. And [King David] mourned over his son a long time.
	}
}
{
	Qere Ketiv.
}

\Verse{38}
{
	\hJ{וְאַבְשָׁלוֹם בָּרַח וַיֵּלֶךְ גְּשׁוּר וַיְהִי־שָׁם שָׁלֹשׁ שָׁנִים׃}
}
{
	\eJ{
		Absalom, who had fled to Geshur, remained there three years.
	}
}
{
}

\Verse{39}
{
	\hJ{וַתְּכַל דָּוִד הַמֶּלֶךְ לָצֵאת אֶל־אַבְשָׁלוֹם כִּי־נִחַם עַל־אַמְנוֹן כִּי־מֵת׃}
}
{
	\eJ{
		And King David was pining away for Absalom, for [the king] had
		gotten over Amnon’s death.
	}
}
{
%	\footnote{
%		Footnote 13:39. David's grief bends toward Absalom while Tamar disappears from the
%		story's speech. That silence is one of the ugliest facts in the chapter.
%	}
}
